\documentclass[12pt,a4paper]{amsart}

\usepackage[ngerman]{babel}
\usepackage{amsmath,amssymb,amsthm}
\usepackage{mathtools}
\usepackage[colorlinks=true,linkcolor=blue,citecolor=blue]{hyperref}
\usepackage{geometry}
\geometry{margin=2.5cm}

% -------------------------------------------------------
% Theorem-Umgebungen
% -------------------------------------------------------
\newtheorem{theorem}{Satz}[section]
\newtheorem{lemma}[theorem]{Lemma}
\newtheorem{korollar}[theorem]{Korollar}
\newtheorem{proposition}[theorem]{Proposition}
\theoremstyle{definition}
\newtheorem{definition}[theorem]{Definition}
\newtheorem{bemerkung}[theorem]{Bemerkung}

% -------------------------------------------------------
% Eigene Makros
% -------------------------------------------------------
\newcommand{\N}{\mathbb{N}}
\newcommand{\Z}{\mathbb{Z}}
\newcommand{\Q}{\mathbb{Q}}
\newcommand{\Fp}{\mathbb{F}_p}
\DeclareMathOperator{\ord}{ord}
\DeclareMathOperator{\ggT}{ggT}

% -------------------------------------------------------
\begin{document}
% -------------------------------------------------------

\title{Der allgemeine Wilson-Satz für endliche abelsche Gruppen}

\author{Michael Fuhrmann}
\address{Specialist-Mathematics-Projekt, Build 65}
\email{specialist-maths@localhost}
\date{\today}

\begin{abstract}
Wir beweisen eine allgemeine Version des Satzes von Wilson, die für jede
endliche abelsche Gruppe $G$ gilt: Das Produkt aller Elemente von $G$ ist
gleich dem neutralen Element $e$, es sei denn, $G$ hat \emph{genau ein}
Element der Ordnung $2$; in diesem Fall ist das Produkt gleich diesem Element.
Der Beweis verwendet nur die Paarungsmethode und elementare Gruppentheorie.
Als Korollar leiten wir den klassischen Satz von Wilson (für $G = (\Z/p\Z)^*$),
den Wilson-Satz für Primzahlpotenzen und die Gaußsche Verallgemeinerung ab:
Für $n \geq 3$ gilt
\[
  \prod_{\substack{a=1\\\ggT(a,n)=1}}^{n} a
  \;\equiv\;
  \begin{cases}
    -1 \pmod{n} & \text{falls } n = 1, 2, 4, p^k \text{ oder } 2p^k
                  \;(p \text{ ungerade Primzahl, } k \geq 1),\\
    \phantom{-}1 \pmod{n} & \text{andernfalls.}
  \end{cases}
\]
\end{abstract}

\maketitle
\tableofcontents

% ================================================================
\section{Einleitung}
% ================================================================

Der klassische Satz von Wilson besagt, dass das Produkt aller Nicht-Null-Elemente
von $\Z/p\Z$ gleich $-1$ ist. Dieses Produkt ist nichts anderes als das Produkt
aller Elemente der multiplikativen Gruppe $G = (\Z/p\Z)^*$.
Der natürliche Rahmen für solche Produktberechnungen ist die Theorie der endlichen
abelschen Gruppen, und das allgemeine Ergebnis ist als \emph{allgemeiner Wilson-Satz}
bekannt.

Die entscheidende Beobachtung ist einfach: In jeder Gruppe heben sich bei der
Multiplikation aller Elemente die nicht-selbst-inversen Elemente paarweise auf.
Nur die selbst-inversen Elemente (diejenigen, die $g^2 = e$ erfüllen) tragen
zum endgültigen Produkt bei. Ob dieser Beitrag $e$ oder etwas Nichttriviales ist,
hängt von der Anzahl der selbst-inversen Elemente ab.

Die Gaußsche Verallgemeinerung des Satzes von Wilson auf zusammengesetzte Moduli
ist dann ein unmittelbares Korollar, sobald wir die Struktur von $(\Z/n\Z)^*$
via dem Chinesischen Restsatz kennen.

% ================================================================
\section{Selbst-inverse Elemente und das Produkt-Lemma}
% ================================================================

\begin{lemma}[Paarungs-Lemma]\label{lem:paarung}
  Sei $(G, \cdot)$ eine endliche abelsche Gruppe mit neutralem Element $e$. Dann gilt:
  \[
    \prod_{g \in G} g \;=\; \prod_{\substack{g \in G \\ g = g^{-1}}} g.
  \]
  Das Produkt aller Elemente von $G$ ist also gleich dem Produkt aller
  selbst-inversen Elemente.
\end{lemma}

\begin{proof}
  Da $G$ abelsch ist, können wir Faktoren frei umsortieren.
  Wir ordnen jedem Element $g \in G$ sein Inverses $g^{-1}$ zu.
  Falls $g \neq g^{-1}$, ist $\{g, g^{-1}\}$ eine zweielementige Teilmenge
  von $G$, und das Paar trägt $g \cdot g^{-1} = e$ zum Produkt bei.
  Da $e$ das Produkt nicht ändert, tragen nur die Elemente mit $g = g^{-1}$
  (d.\,h.\ $g^2 = e$) nicht-trivial bei.
  Daher gilt $\prod_{g \in G} g = \prod_{g^2=e} g$.
\end{proof}

\begin{lemma}[Untergruppe der Involutionen]\label{lem:untergruppe}
  Sei $(G, \cdot)$ eine endliche abelsche Gruppe. Die Menge
  $S = \{g \in G : g^2 = e\}$ ist eine Untergruppe von $G$. Außerdem ist
  $S$ eine elementar-abelsche $2$-Gruppe: entweder $S = \{e\}$ oder
  $S \cong (\Z/2\Z)^k$ für ein $k \geq 1$. Insbesondere gilt $|S| = 2^k$
  für ein $k \geq 0$.
\end{lemma}

\begin{proof}
  \textbf{$S$ ist eine Untergruppe:} $e^2 = e$, also $e \in S$.
  Für $g, h \in S$ gilt $(gh)^2 = g^2 h^2 = e \cdot e = e$ (wegen Kommutativität),
  also $gh \in S$. Falls $g \in S$, dann $g^{-1} = g$ (da $g^2 = e$ impliziert
  $g = g^{-1}$), also $g^{-1} \in S$. Damit ist $S \leq G$.

  \textbf{$S$ ist elementar-abelsch:} Jedes Nicht-Einheitselement von $S$ hat
  Ordnung genau $2$. Daher ist $S$ eine abelsche Gruppe, in der jedes Element
  Ordnung $1$ oder $2$ hat. Nach dem Klassifikationssatz endlicher abelscher
  Gruppen ist $S$ ein direktes Produkt von Kopien von $\Z/2\Z$:
  $S \cong (\Z/2\Z)^k$. Insbesondere $|S| = 2^k$.
\end{proof}

% ================================================================
\section{Der allgemeine Wilson-Satz}
% ================================================================

\begin{theorem}[Allgemeiner Wilson-Satz]\label{satz:allgemeiner_wilson}
  Sei $(G, \cdot)$ eine endliche abelsche Gruppe. Dann gilt:
  \[
    \prod_{g \in G} g
    \;=\;
    \begin{cases}
      e   & \text{falls die Anzahl der Elemente der Ordnung } 2 \text{ in } G
             \text{ nicht gleich } 1 \text{ ist}, \\
      \tau & \text{falls } G \text{ genau ein Element } \tau \text{ der Ordnung } 2 \text{ hat.}
    \end{cases}
  \]
\end{theorem}

\begin{proof}
  Nach Lemma~\ref{lem:paarung} gilt $\prod_{g \in G} g = \prod_{g \in S} g$,
  wobei $S = \{g \in G : g^2 = e\}$.

  Nach Lemma~\ref{lem:untergruppe} gilt $S \cong (\Z/2\Z)^k$ für ein $k \geq 0$.

  \textbf{Fall $k = 0$:} $S = \{e\}$, also $\prod_{g \in S} g = e$.

  \textbf{Fall $k = 1$:} $S = \{e, \tau\}$, wobei $\tau$ die Ordnung $2$ hat.
  Dann $\prod_{g \in S} g = e \cdot \tau = \tau$.

  \textbf{Fall $k \geq 2$:} $S \cong (\Z/2\Z)^k$ mit $k \geq 2$.
  Wir behaupten $\prod_{g \in S} g = e$.

  Sei $\tau \in S \setminus \{e\}$ ein beliebiges Element der Ordnung $2$.
  Die Abbildung $\varphi: S \to S$, $s \mapsto \tau s$, ist eine Bijektion
  (da $\tau \in S$ und $S$ abgeschlossen unter Multiplikation ist).
  Die Fixpunkte von $\varphi$ sind die $s$ mit $\tau s = s$, also $\tau = e$,
  was unmöglich ist. Daher hat $\varphi$ keine Fixpunkte, und $S$ zerfällt in
  Paare $\{s, \tau s\}$.

  Das Produkt eines solchen Paares ist $s \cdot (\tau s) = \tau s^2 = \tau e = \tau$,
  da $s^2 = e$ für alle $s \in S$.

  $S$ zerfällt in $|S|/2 = 2^{k-1}$ solche Paare, jedes mit Paarprodukt $\tau$.
  Daher:
  \[
    \prod_{g \in S} g = \tau^{2^{k-1}}.
  \]
  Da $\tau^2 = e$ und $k \geq 2$, ist $2^{k-1}$ gerade, also
  $\tau^{2^{k-1}} = (\tau^2)^{2^{k-2}} = e^{2^{k-2}} = e$.

  In allen Fällen ist die Formel bewiesen. \qedhere
\end{proof}

\begin{bemerkung}
  Die Anzahl der Elemente der Ordnung $2$ in einer endlichen abelschen Gruppe
  ist immer von der Form $2^k - 1$ für ein $k \geq 0$ (nur Elemente der Ordnung
  genau $2$, nicht $1$ gezählt): für $k=0$ gibt es $0$ solche Elemente,
  für $k=1$ genau $1$, für $k=2$ genau $3$, usw. Insbesondere ist die Anzahl
  \emph{nie} gleich $2$, daher ist der einzige Fall, in dem das Produkt
  nicht-trivial ist, genau dann wenn es genau ein Element der Ordnung $2$ gibt.
\end{bemerkung}

% ================================================================
\section{Klassischer Wilson als Spezialfall}
% ================================================================

\begin{korollar}[Klassischer Wilson-Satz]\label{kor:klassischer_wilson}
  Für eine ungerade Primzahl $p$ gilt: $(p-1)! \equiv -1 \pmod{p}$.
\end{korollar}

\begin{proof}
  Setze $G = (\Z/p\Z)^*$. Diese Gruppe ist zyklisch von gerader Ordnung $p-1$.
  Eine zyklische Gruppe gerader Ordnung hat genau ein Element der Ordnung $2$.
  Für $G = (\Z/p\Z)^*$ ist dieses Element $p-1 \equiv -1 \pmod{p}$.
  Nach Satz~\ref{satz:allgemeiner_wilson} (Fall $k=1$) ist das Produkt aller
  Elemente gleich $-1 \pmod{p}$, was $(p-1)!$ ergibt.
\end{proof}

% ================================================================
\section{Die Gaußsche Verallgemeinerung}
% ================================================================

Wir charakterisieren nun alle $n$, für die das Produkt der Einheiten modulo $n$
gleich $-1$ ist.

\begin{theorem}[Gaußsche Verallgemeinerung von Wilson]\label{satz:gauss}
  Für $n \geq 1$ gilt:
  \[
    \prod_{\substack{a=1\\\ggT(a,n)=1}}^{n} a
    \;\equiv\;
    \begin{cases}
      -1 \pmod{n} & \text{falls } n \in \{1, 2, 4\}
                    \text{ oder } n = p^k \text{ oder } n = 2p^k
                    \;(p \text{ ungerade Primzahl, } k \geq 1),\\
      \phantom{-}1 \pmod{n} & \text{andernfalls.}
    \end{cases}
  \]
\end{theorem}

\begin{proof}
  Das Produkt $\prod_{\ggT(a,n)=1} a$ ist das Produkt aller Elemente von
  $G_n = (\Z/n\Z)^*$. Nach Satz~\ref{satz:allgemeiner_wilson} ist dieses
  Produkt genau dann gleich $-1$, wenn $G_n$ genau ein Element der Ordnung $2$ hat.

  Wir bestimmen, wann $G_n$ genau ein Element der Ordnung $2$ hat, mit Hilfe des
  Struktursatzes für $(\Z/n\Z)^*$.

  \textbf{Schritt 1: Struktur von $(\Z/n\Z)^*$.}
  Nach dem Chinesischen Restsatz gilt für $n = \prod_{i} p_i^{k_i}$:
  \[
    (\Z/n\Z)^* \;\cong\; \prod_i (\Z/p_i^{k_i}\Z)^*.
  \]
  Die Struktur der einzelnen Faktoren ist:
  \begin{itemize}
    \item $(\Z/2\Z)^* = \{1\}$ (trivial),
    \item $(\Z/4\Z)^* \cong \Z_2$,
    \item $(\Z/2^k\Z)^* \cong \Z_2 \times \Z_{2^{k-2}}$ für $k \geq 3$,
    \item $(\Z/p^k\Z)^* \cong \Z_{p^{k-1}(p-1)}$ für ungerades $p$ (zyklisch).
  \end{itemize}

  \textbf{Schritt 2: Elemente der Ordnung 2.}
  Die Anzahl der Elemente der Ordnung höchstens $2$ in einem Produkt
  $H_1 \times \cdots \times H_r$ ist $\prod_i |\{h \in H_i : h^2 = e\}|$.

  Für jeden Faktor:
  \begin{itemize}
    \item $(\Z/2\Z)^*$: $1$ Element der Ordnung $\leq 2$ (nämlich $e$).
    \item $(\Z/4\Z)^* \cong \Z_2$: $2$ Elemente der Ordnung $\leq 2$.
    \item $(\Z/2^k\Z)^*$ für $k \geq 3$: $4$ Elemente der Ordnung $\leq 2$.
    \item $(\Z/p^k\Z)^*$ für ungerades $p$: zyklisch von gerader Ordnung,
      also $2$ Elemente der Ordnung $\leq 2$ (nämlich $e$ und $-1$).
  \end{itemize}

  Die Gesamtanzahl der Elemente der Ordnung $\leq 2$ in $G_n$ ist das Produkt
  dieser Zählungen über alle Primzahlpotenz-Faktoren.

  \textbf{Schritt 3: Wann ist die Anzahl genau 1?}
  Wir brauchen die Gesamtanzahl der Elemente der Ordnung $\leq 2$ gleich $2$
  (eins für $e$, eins für $\tau$). Das Produkt muss also $2$ ergeben.
  Daher trägt genau ein Faktor $2$ bei, alle anderen $1$.

  Ein Faktor trägt $1$ bei genau dann, wenn die entsprechende Gruppenkomponente
  trivial ist (kein Element der Ordnung $2$ außer $e$): das ist $(\Z/2\Z)^*$.
  Ein Faktor trägt $2$ bei für $(\Z/4\Z)^*$ oder $(\Z/p^k\Z)^*$ mit ungeradem $p$.
  Ein Faktor trägt $\geq 4$ bei für $(\Z/2^k\Z)^*$ mit $k \geq 3$.

  Damit das Produkt gleich $2$ ist:
  \begin{itemize}
    \item Es gibt keinen Faktor mit Beitrag $\geq 4$: keine $2^k$-Komponente
      mit $k \geq 3$.
    \item Genau ein Faktor mit Beitrag $2$: entweder genau ein ungerader
      Primzahlpotenzen-Faktor (und kein $4$-Faktor), oder genau ein $4$-Faktor
      (und keine ungeraden Primzahlpotenz-Faktoren).
  \end{itemize}

  Die $n$-Werte, die diese Bedingung erfüllen, sind:
  \begin{itemize}
    \item $n = 1$: trivialerweise gilt $\prod = 1 \equiv -1 \pmod{1}$.
    \item $n = 2$: $G_2 = \{1\}$, Produkt $1 \equiv -1 \pmod{2}$.
    \item $n = 4$: $G_4 \cong \Z_2$, eindeutiges Element der Ordnung $2$ ist
      $3 \equiv -1 \pmod{4}$. \checkmark
    \item $n = p^k$ ($p$ ungerade Primzahl, $k \geq 1$): $G_{p^k}$ ist zyklisch
      von gerader Ordnung, eindeutiges Element der Ordnung $2$ ist $-1$. \checkmark
    \item $n = 2p^k$ ($p$ ungerade Primzahl, $k \geq 1$): $G_{2p^k} \cong
      \{1\} \times G_{p^k}$ da $(\Z/2\Z)^*$ trivial ist. Also $G_{2p^k} \cong
      G_{p^k}$, zyklisch mit eindeutigem Element der Ordnung $2$, nämlich $-1$. \checkmark
  \end{itemize}

  Alle anderen $n$ haben entweder $0$ oder $\geq 3$ Elemente der Ordnung $2$
  in $G_n$, was das Produkt $e = 1 \pmod{n}$ ergibt. \qedhere
\end{proof}

\begin{korollar}[Produkt $\equiv 1$ für die meisten $n$]\label{kor:produkt_eins}
  Für $n \geq 3$ mit $n \neq 4$, $n \neq p^k$, $n \neq 2p^k$ gilt:
  \[
    \prod_{\substack{a=1\\\ggT(a,n)=1}}^{n} a \;\equiv\; 1 \pmod{n}.
  \]
  Insbesondere ist für $n = 8$ das Produkt aller ungeraden Reste modulo $8$
  gleich $1 \cdot 3 \cdot 5 \cdot 7 = 105 \equiv 1 \pmod{8}$.
\end{korollar}

\begin{proof}
  Dies ist die Komplementaussage von Satz~\ref{satz:gauss}. Für $n = 8$:
  $G_8 = (\Z/8\Z)^*$ hat vier Elemente der Ordnung $\leq 2$ (nämlich $1,3,5,7$),
  mehr als eines der Ordnung $2$, also ist das Produkt $1$.
  Direkter Check: $1 \cdot 3 \cdot 5 \cdot 7 = 105 = 13 \cdot 8 + 1 \equiv 1 \pmod{8}$. \checkmark
\end{proof}

% ================================================================
\section{Schlussfolgerung}
% ================================================================

Der allgemeine Wilson-Satz vereint alle spezifischen Fälle in einer einzigen
gruppentheoretischen Aussage: Das Produkt aller Elemente einer endlichen
abelschen Gruppe $G$ ist nicht-trivial (gleich $\tau$) genau dann, wenn $G$
genau eine Involution $\tau$ (Element der Ordnung $2$) hat. Nach dem
Struktursatz für $(\Z/n\Z)^*$ tritt dies genau für $n = 4, p^k, 2p^k$ auf.

Die Beweismethode – Paarung von Elementen mit ihren Inversen – ist elementar
aber wirkungsvoll. Sie funktioniert in jeder Gruppe (nicht nur in abelschen),
aber die abschließende Identifikation der ungepaaraten Elemente erfordert
Kommutativität, um sicherzustellen, dass $|S|$ immer eine Zweierpotenz ist.

\begin{thebibliography}{9}

\bibitem{Hardy1979}
G.~H.~Hardy und E.~M.~Wright.
\newblock \emph{Einführung in die Zahlentheorie} (An Introduction to the Theory of Numbers), 5.~Aufl.
\newblock Oxford University Press, 1979.

\bibitem{Ireland1990}
K.~Ireland und M.~Rosen.
\newblock \emph{A Classical Introduction to Modern Number Theory}, 2.~Aufl.
\newblock Springer, New York, 1990.

\bibitem{Lang2002}
S.~Lang.
\newblock \emph{Algebra}, revidierte 3.~Aufl.
\newblock Springer, New York, 2002.

\bibitem{Dummit2003}
D.~S.~Dummit und R.~M.~Foote.
\newblock \emph{Abstract Algebra}, 3.~Aufl.
\newblock Wiley, Hoboken, 2003.

\end{thebibliography}

\end{document}
